\documentclass[10pt]{article}
% NOTE: The official ICML/ICLR/NeurIPS style files were not available offline in the
% sandbox (no network), so this uses the generic article class with a similar two-column-free
% layout. Swap in the venue's official style file before submission.
\usepackage[margin=1in]{geometry}
\usepackage{amsmath,amssymb,booktabs}
\usepackage[hidelinks]{hyperref}
\title{Simple Statistics Predict Random-Walk Cover Times on Small Graphs}
\author{EvoScientist (anonymous draft)}
\date{}
\begin{document}
\maketitle

\begin{abstract}
The expected cover time $C(G)$ of a simple random walk on a connected graph is classically
bounded by $2m(n-1)$, which is loose for small graphs. We ask whether a regression on a few
cheap statistics (vertex count $n$, edge count $m$, maximum degree $\Delta$, diameter $D$) can
predict $C(G)$ more accurately for connected graphs with $n\le 30$, and where such a predictor
fails. We generate 168 graphs from 11 families, estimate $C(G)$ by vectorised Monte Carlo
(1000 walks per start vertex, worst start), validate the estimator against exact linear-system
values for the 68 graphs with $n\le 10$, and evaluate log-linear predictors with
leave-one-family-out (LOFO) cross-validation. The classic bound overestimates $C(G)$ by a median
factor of about $3.6$ (median relative error $0.53$). A three-feature log-linear model
$\log C\approx a+b\log n+c\log m+d\log D$ reduces LOFO median relative error to $0.19$
(95\% bootstrap CI $[0.16,0.24]$), and adding $\log\Delta$ to give four features reduces it to $0.16$.
Error is strongly family-dependent: the four-feature model is accurate on paths, cycles, random trees and
sparse random graphs (median error $0.04$--$0.09$), but the three-feature model is poor on cycles ($0.58$),
and both models are poor on grids, lollipops, $G(n,p)$ and the complete graph
(median error $0.5$--$0.7$), where cover time is governed by bottlenecks or density
effects that these global statistics do not capture. All conclusions are specific to this small-graph
synthetic benchmark.
\end{abstract}

\section{Introduction}
For a random walk on a connected undirected graph, the cover time is the expected number of
steps needed to visit every vertex, maximised over start vertices. Aleliunas et al.\ (1979)
showed that it is at most $2m(n-1)$, and later work (e.g.\ Feige, 1995; Lovász, 1993) gives
sharper asymptotic characterisations in terms of hitting times. These results are asymptotic or
worst-case over graph classes, so for a specific graph with $n\le 30$ they can be far from the
true value. A practitioner who needs a quick estimate of cover time for a small network would
benefit from a predictor that uses only a handful of easily computed statistics and whose error
is known per graph type. We provide such an empirical predictor and measure its reliability.

\textbf{Research question.} For connected graphs with at most 30 vertices, can a few cheap
statistics ($n$, $m$, mean/max degree, diameter) predict $C(G)$ more accurately than
$2m(n-1)$, and how does the error depend on graph type?

\textbf{Contributions.} (i) A reproducible Monte Carlo benchmark of $C(G)$ over 168 graphs,
validated against exact values for $n\le 10$. (ii) A comparison of the classic bound with
simple log-linear predictors under family-held-out evaluation. (iii) A family-level error
analysis that identifies where the predictors fail.

\section{Method}
\textbf{Estimand.} $C(G)=\max_{s} \mathbb{E}_s[\tau_{\text{cover}}]$. For $n\le 8$ all starts are
simulated; for larger $n$ we simulate four candidate starts (both ends of a double-sweep diameter
path, a maximum-degree vertex, and one random vertex) and take the largest mean. Selecting the
maximum over starts introduces a small upward bias; we keep it because it targets the definition.

\textbf{Monte Carlo.} $R=1000$ independent walks per start are simulated in vectorised NumPy until
every vertex is visited. The standard error of the mean is reported. The median relative SE is
$0.017$ (max $0.032$).

\textbf{Exact validation.} For $n\le 10$ we solve the linear system for the expected remaining
cover time over states $(v,S)$, $v\in S$, with sparse LU, and take the maximum over starts.
Across 68 graphs, mean $|C_{\text{MC}}/C_{\text{exact}}-1|=0.013$ (seed 0) and $0.015$ (seed 1);
the largest $|z|$ is $2.1$ and $2.6$, respectively, and 98.5\% (seed 0) / 95.6\% (seed 1) of
estimates fall within two standard errors of the exact value.

\textbf{Graphs.} Eight deterministic families (path, cycle, star, complete, $\sqrt n\times\sqrt n$ grid,
lollipop, barbell, binary tree) and three random families (random recursive tree, $G(n,p)$ with
$p=2\ln n/n+0.05$ conditioned on connectivity, and random tree plus $\lfloor n/4\rfloor$ extra edges).
Sizes $n\in\{4,6,8,10,12,16,20,24,30\}$; random families get 2 or 4 instances per size.

\textbf{Predictors.} All are OLS fits of $\log \hat C$ on:
\begin{itemize}
\item \texttt{bound\_scale}: $\log(2m(n-1))$ with one fitted scale and intercept;
\item \texttt{n\_m}: $\log n$, $\log m$;
\item \texttt{n\_m\_D}: $\log n$, $\log m$, $\log D$;
\item \texttt{full}: additionally $\log \Delta$.
\end{itemize}
Evaluation uses leave-one-family-out CV (fit on 10 families, predict the held-out one), so each
prediction is made for a family not seen in training. Uncertainty over graphs is a 2000-sample
bootstrap of the median LOFO error. Relative error is $|\hat C_{\text{pred}}-\hat C_{\text{MC}}|/\hat C_{\text{MC}}$.

\section{Experiments}
Table~\ref{tab:main} gives the main results (seed 0; seed 1 in brackets). Table~\ref{tab:fam}
gives per-family LOFO median errors for seed 0.

\begin{table}[h]
\centering
\caption{Predictor accuracy over 168 graphs. LOFO = leave-one-family-out. CI = 95\% bootstrap interval for the LOFO median (seed 0).}
\label{tab:main}
\begin{tabular}{lccc}
\toprule
Predictor & In-sample median err. & LOFO median err. & LOFO 90th pct. \\
\midrule
Classic bound $2m(n-1)$ (scaled) & 0.48 [0.50] & 0.53 [0.56] (CI 0.50--0.61) & 2.81 [2.83] \\
$\log n,\log m$ & 0.20 [0.20] & 0.23 [0.26] (CI 0.19--0.32) & 0.83 [0.84] \\
$\log n,\log m,\log D$ & 0.14 [0.14] & 0.19 [0.18] (CI 0.16--0.24) & 0.76 [0.73] \\
$+\log\Delta$ & 0.13 & 0.16 (CI 0.11--0.20) & 0.74 \\
\bottomrule
\end{tabular}
\end{table}

\begin{table}[h]
\centering
\caption{Per-family LOFO median relative error (seed 0), for the three-feature model (\texttt{n\_m\_D}), the four-feature model (\texttt{full}), and the classic bound (scaled).}
\label{tab:fam}
\begin{tabular}{lccc}
\toprule
Family & n\_m\_D & full & bound (scaled) \\
\midrule
sparse random & 0.09 & 0.09 & 0.38 \\
random tree & 0.13 & 0.07 & 0.52 \\
path & 0.13 & 0.04 & 0.62 \\
binary tree & 0.13 & 0.25 & 0.47 \\
complete & 0.20 & 0.48 & 5.26 \\
barbell & 0.27 & 0.31 & 0.59 \\
star & 0.52 & 0.16 & 0.12 \\
$G(n,p)$ & 0.54 & 0.66 & 1.90 \\
cycle & 0.58 & 0.06 & 0.41 \\
lollipop & 0.62 & 0.53 & 0.68 \\
grid & 0.71 & 0.67 & 0.38 \\
\bottomrule
\end{tabular}
\end{table}

Fitted exponents for \texttt{n\_m\_D} on all data are $b_n=1.35$, $b_m=0.07$, $b_D=0.78$
(seed 1: $1.33$, $0.08$, $0.77$). The edge exponent is near zero once $n$ and $D$ are known,
so $C(G)$ is governed mainly by size and diameter in this benchmark. The classic bound
has the right order of growth in $m$ and $n$ but overestimates by a graph-dependent factor:
the ratio $C/[2m(n-1)]$ has median $0.27$ and ranges from $0.005$ (complete) to $0.64$ (path).
Residuals of \texttt{n\_m\_D} correlate weakly with $\Delta/n$ ($r=0.30$), and not with
$\log(m/n)$ ($r=-0.05$) or $D/n$ ($r=0.03$), which is consistent with the four-feature
model's gain coming from degree heterogeneity.

\section{Discussion and limitations}
A three-feature log-linear model roughly halves the error of the classic bound on
held-out families, and this holds on both seeds. The four-feature model is reliable on paths, cycles, random trees and sparse random graphs, where
the cover time is driven mainly by size and diameter (the three-feature model is not reliable on cycles).
It is unreliable on grids, lollipops and barbells, which contain bottlenecks that the global
statistics do not see, and on dense graphs ($G(n,p)$, complete), where the within-family
relationship differs from the sparse families and LOFO extrapolation fails. The 90th-percentile
errors (about $0.75$) show that individual predictions can be off by a large factor, so the predictor
should be used only as a rough estimate with an error band.

Limitations: (1) the benchmark is synthetic and small ($n\le 30$), and the four-feature model's
coefficients are fitted to these families only; (2) the target is the maximum over starts, which
for $n>8$ is taken over four candidate starts, so it may slightly overstate the true maximum;
(3) only two Monte Carlo seeds of graph generation were run, with one simulation seed per
graph; (4) the families were chosen by us, and LOFO errors depend on this choice; (5) we did not
compare to other baseline predictors (e.g.\ hitting-time-based Matthews bounds or spectral
quantities), so claims of superiority are relative only to the classic bound and to
log-linear fits on these statistics. Exact validation was limited to $n\le 10$.

\section{Conclusion}
On small connected graphs, a log-linear predictor in $n$, $m$ and diameter (optionally
$\Delta$) gives substantially better cover-time estimates than $2m(n-1)$, with median
held-out relative error of about 0.16--0.19. Its error is strongly family-dependent: it is
accurate on paths, cycles, random trees and sparse random graphs under the four-feature model,
and not reliable on bottleneck or dense graphs.

\section*{Reproducibility}
Code: \texttt{cover\_experiment.py} (simulation and exact solver), \texttt{analyze.py} (fits and CV).
Commands: \texttt{python3 cover\_experiment.py --seed 0 --R 1000 --out results.csv};
\texttt{python3 analyze.py results.csv > analysis.txt}; the same with \texttt{--seed 1} for the robustness run.
Seeds: graph generation uses \texttt{seed}; Monte Carlo uses \texttt{seed+1000}; bootstrap uses 12345.
Software: Python 3.11, NumPy 2.4, SciPy 1.17, pandas. Runtime: about 1--2 minutes per seed on CPU.

\begin{thebibliography}{9}
\bibitem{aleliunas} R.~Aleliunas, R.~M.~Karp, R.~J.~Lipton, L.~Lov\'asz, C.~Rackoff.
Random walks, universal traversal sequences, and the complexity of maze problems. FOCS, 1979.
\bibitem{feige} U.~Feige. A tight upper bound on the cover time for random walks on graphs.
Random Structures \& Algorithms, 1995.
\bibitem{lovasz} L.~Lov\'asz. Random walks on graphs: a survey. Combinatorics, Paul Erd\H{o}s is Eighty, 1993.
\end{thebibliography}
\end{document}
